% ===========================================================================
% ABSTRACT
% File: sections/00_abstract.tex
% ===========================================================================

\begin{abstract}
\noindent
Land-use competition between food production and renewable energy deployment
is among the most pressing resource conflicts of the energy transition,
particularly in semi-arid regions where solar irradiance is high but
agricultural land is limited. This study develops and applies an integrated
optimization framework for bifacial agrivoltaic (APV) and on-farm anaerobic
digestion (AD) systems across four distinct climate zones: Konya, Turkey
(K\"{o}ppen BSk, cold semi-arid), Almer\'{i}a, Spain (BSh, hot semi-arid),
Ouagadougou, Burkina Faso (BSh, tropical semi-arid), and Freiburg,
Germany (Cfb, oceanic temperate).

Hourly PVGIS-SARAH3 Typical Meteorological Year data drive a coupled
six-model simulation pipeline covering bifacial PV power, tomato crop
yield via PAR integral, evapotranspiration reduction, anaerobic digestion
with Arrhenius temperature correction, and digester thermal management.
Six design variables --- panel tilt ($\beta$), row spacing
($d_{\text{row}}$), module height ($H_m$), digester volume
($V_{\text{dig}}$), hydraulic retention time (HRT), and PV heating
fraction ($f_{\text{PV,heat}}$) --- are jointly optimized using
Differential Evolution to maximize a novel three-component extended Land
Equivalent Ratio:
\[
    eLER = LER_{\text{crop}} + LER_{\text{PV}} + LER_{\text{biogas}}
\]

Optimized eLER values reach 1.653 (Konya), 1.659 (Almer\'{i}a), 1.725
(Ouagadougou), and 1.758 (Freiburg). All sites exceed the two-component
agrivoltaic benchmark when the biogas term is included, and
$LER_{\text{2C}}$ alone ranges from 1.031 to 1.089, confirming that the
APV system outperforms two separate monocultures at every site. A
counter-intuitive climate-gradient result emerges: Freiburg, with the
lowest annual solar resource (\SI{1150}{kWh.m^{-2}}), achieves the
highest eLER (1.758), exceeding Konya (\SI{1650}{kWh.m^{-2}}, eLER =
1.653). This inversion is driven by PAR saturation: at high irradiance,
the open-field PAR integral exceeds the tomato light saturation point
(\SI{174}{W.m^{-2}}), inflating the $LER_{\text{crop}}$ denominator and
penalizing high-irradiance sites.

The optimizer converges to the physical constraint boundary
($d_{\text{row}} = \SI{3.03}{m}$, $H_m = \SI{1.50}{m}$) at all four
sites, a structural boundary-convergence result that quantifies the eLER
penalty of any additional mechanical access requirement. A 4E analysis
(Energetic, Economic, Energo-economic, Environmental) reveals annual PV
yields of 1{,}354--1{,}583~MWh\,ha$^{-1}$, NPV values of
\SIrange{169}{981}{k\$}, and IRR of \SIrange{11.8}{31.1}{\percent} under
the baseline scenario (S0). The carbon-credit scenario (S6) is the
economically dominant strategy at all sites, raising IRR to
\SIrange{13.5}{33.1}{\percent}. Annual CO$_2$ avoidance ranges from 251
(Almer\'{i}a) to 951 (Ouagadougou)\,tCO$_2$\,ha$^{-1}$. Growing-season
water savings of \SIrange{152}{383}{mm} represent
\SIrange{15.5}{25.8}{\percent} of seasonal ET$_0$. Crop yield model
validation (MAPE\,=\,2.4\% in the operational shading range) confirms
that all reported results are conservative lower bounds.

\bigskip
\noindent\textbf{Keywords:} agrivoltaics; anaerobic digestion; extended
land equivalent ratio; food--energy--water nexus; differential evolution;
semi-arid climate; PAR saturation; 4E analysis; biomethane; carbon credit
\end{abstract}
